\section*{Capítulo \ref{ch:b3:surfaces-catalysis} --- Superfícies, adsorção e catálise heterogênea}

\begin{solution}{exo:b3:surfaces-catalysis:1}
$K = \theta/((1 - \theta)p) = 0.40/(0.60 \times 2.0) = \qty{0.33}{kPa^{-1}}$; a \qty{10}{kPa}, $\theta = 3.33/4.33 = 0.77$.
\end{solution}

\begin{solution}{exo:b3:surfaces-catalysis:2}
$-20$: \omterm{def:b3:surfaces-catalysis:physi-chemi}{fisissorção} (a ordem de grandeza de uma entalpia de condensação); $-150$: \omterm{def:b3:surfaces-catalysis:physi-chemi}{quimissorção},
a única que pode ser dissociativa.
\end{solution}

\begin{solution}{exo:b3:surfaces-catalysis:3}
(100): um sítio de topo, dois de ponte (quatro ligações compartilhadas por dois átomos
cada), uma cavidade quaternária (quatro cavidades compartilhadas por quatro átomos).
(111): um de topo, três de ponte, duas cavidades ternárias (seis triângulos em torno de
cada átomo, cada um compartilhado por três).
\end{solution}

\begin{solution}{exo:b3:surfaces-catalysis:4}
$\num{3.0e-6}/\num{1.5e-5} = \qty{0.20}{s^{-1}}$.
\end{solution}

\begin{solution}{exo:b3:surfaces-catalysis:5}
$\theta/(1 - \theta) = (Kp)^{1/2}$: $0.111$ a \qty{1.0}{kPa}, o dobro a \qty{4.0}{kPa}: $0.222$, logo $\theta = 0.18$.
\end{solution}

\begin{solution}{exo:b3:surfaces-catalysis:6}
Denominador $1 + 10 + 2.5 = 13.5$: $\theta_{\ce{CO}} = 0.74$, $\theta_{\ce{O2}} = 0.19$, sítios livres 0.07. A
velocidade de Langmuir--Hinshelwood, proporcional a $\theta_{\ce{CO}}\theta_{\ce{O2}}$, é limitada pela
escassez de oxigênio em uma superfície tomada pelo CO.
\end{solution}

\begin{solution}{exo:b3:surfaces-catalysis:7}
$\Delta_{\mathrm{ads}}H = -R\ln(8.0/1.0)/(1/300 - 1/350) = -8.314 \times 2.079/\num{4.76e-4} = \qty{-36}{kJ/mol}$.
\end{solution}

\begin{solution}{exo:b3:surfaces-catalysis:8}
$n_{\mathrm m} = 120/22\,414 = \qty{5.35e-3}{mol}$; área $= \num{5.35e-3} \times \num{6.022e23} \times \num{0.162e-18} =
\qty{522}{m^2.g^{-1}}$.
\end{solution}

\begin{solution}{exo:b3:surfaces-catalysis:9}
Langmuir--Hinshelwood: os dois reagentes precisam estar adsorvidos em sítios vizinhos;
o CO se liga fortemente e, a alta pressão, não deixa lugar para o oxigênio.
\end{solution}

\begin{solution}{exo:b3:surfaces-catalysis:10}
Baixo recobrimento: $100 - 60 = \qty{40}{kJ/mol}$; alto recobrimento: \qty{100}{kJ/mol}. O gráfico de $\ln r$
em função de $1/T$ é íngreme (inclinação $-100/R$) a baixa temperatura, onde a superfície
está cheia, e mais plano (inclinação $-40/R$) a alta temperatura, onde ela se esvazia:
uma curva que se dobra entre os dois regimes.
\end{solution}

\begin{solution}{exo:b3:surfaces-catalysis:11}
$4 \times 0.10 = 0.40$ dos sítios estão bloqueados: resta 60\,\% da atividade para um
sítio, cerca de $0.60^2 = 36\,\%$ para dois sítios livres adjacentes.
\end{solution}

\begin{solution}{exo:b3:surfaces-catalysis:12}
O do meio (\qty{-250}{kJ/mol}): o primeiro liga o nitrogênio fraco demais para romper o
\ce{N2} com facilidade, o terceiro forte demais, de modo que os átomos de nitrogênio
permanecem na superfície e a bloqueiam.
\end{solution}

\begin{solution}{pb:b3:surfaces-catalysis:1}
\textbf{1.} O nitrogênio se condensa perto de \qty{77}{K} à pressão atmosférica: alcança-se
toda a faixa de pressões relativas, e formam-se multicamadas.
\textbf{2.} Abaixo de 0.05, a heterogeneidade da superfície, e acima de 0.30, a
condensação nos poros violam as hipóteses do BET.
\textbf{3.} \num{1.278e-3}, \num{2.410e-3}, \num{3.453e-3}, \num{4.621e-3}, \num{5.659e-3}, \num{6.813e-3} \unit{g.cm^{-3}}.
\textbf{4.} $v_{\mathrm m} = 1/(0.02205 + 0.000180) = \qty{45.0}{cm^3.g^{-1}}$.
\textbf{5.} $c = 1 + 0.02205/0.000180 = 124$: a primeira camada se liga muito mais
fortemente que o líquido.
\textbf{6.} Ele é minúsculo em comparação com a inclinação vezes $x$, de modo que seu erro
relativo é grande; mas $v_{\mathrm m}$ depende da soma da inclinação e do intercepto, dominada
pela inclinação.
\textbf{7.} $45.0/22\,414 = \qty{2.01e-3}{mol.g^{-1}}$.
\textbf{8.} $\num{2.01e-3} \times \num{6.022e23} \times \num{0.162e-18} = \qty{196}{m^2.g^{-1}}$.
\textbf{9.} O suporte de alumina: a platina exposta (Parte III) tem cerca de \qty{0.9}{m^2.g^{-1}}.
\textbf{10.} Poros de largura molecular se enchem completamente a pressão muito baixa,
e não camada por camada, de modo que nenhuma \omterm{def:b3:surfaces-catalysis:coverage}{monocamada} pode ser identificada.
\textbf{11.} $2 \times 0.205/22\,414 = \qty{1.83e-5}{mol.g^{-1}}$.
\textbf{12.} $0.0100/195.08 = \qty{5.13e-5}{mol.g^{-1}}$.
\textbf{13.} $D = 1.83/5.13 = 0.357$.
\textbf{14.} $a^3/4 = (0.3923)^3/4 = \qty{0.01509}{nm^3}$.
\textbf{15.} $\frac{\sqrt3}{4}(0.3923)^2 = \qty{0.0666}{nm^2}$, a face mais densa; (100) e (110) dão 0.0770 e
\qty{0.1088}{nm^2}, e a média das três densidades de sítios dá \qty{0.0807}{nm^2}.
\textbf{16.} Uma esfera contém $\pi d^3/6v_{\mathrm{at}}$ átomos, dos quais $\pi d^2/a_{\mathrm{at}}$ estão na superfície:
$D = 6v_{\mathrm{at}}/(a_{\mathrm{at}}d)$.
\textbf{17.} $d = 6 \times 0.01509/(0.0807 \times 0.357) = \qty{3.1}{nm}$.
\textbf{18.} Um H por Pt de superfície; partículas esféricas de um único tamanho (o
resultado é uma média ponderada pela superfície); toda a platina reduzida e acessível;
nenhum hidrogênio migrando para o suporte.
\textbf{19.} $\num{2.0e-5}/\num{1.83e-5} = \qty{1.1}{s^{-1}}$.
\textbf{20.} $\num{2.0e-5}/0.0100 = \qty{2.0e-3}{mol.g^{-1}.s^{-1}}$ por grama de platina.
\textbf{21.} $D = 1.12/10 = 0.112$ em vez de 0.357: um fator de 3.2.
\textbf{22.} Que a reação é sensível à estrutura: seus \omterm{def:b3:complex-kinetics:enzyme}{sítios ativos} (arestas, vértices,
faces particulares) não são uma fração constante da superfície.
\textbf{23.} Ela conta os eventos por \omterm{def:b3:complex-kinetics:enzyme}{sítio ativo}, eliminando o efeito da quantidade de
metal exposta.
\textbf{24.} \textbf{Diâmetro médio das partículas de platina $\approx \qty{3.1}{nm}$.}
\end{solution}
